\documentclass[10pt,a4paper]{amsart}
\usepackage{amsmath,amssymb,amsthm,graphicx}
\usepackage{verbatim}
\newenvironment{reference}{\comment}{\endcomment}
%\newenvironment{reference}{}{}
\newenvironment{slogan}{\comment}{\endcomment}
\newenvironment{history}{\comment}{\endcomment}

% For commutative diagrams we use Xy-pic
\usepackage[all]{xy}

% We use 2cell for 2-commutative diagrams.
\xyoption{2cell}
% Theorem environments.
%
\theoremstyle{plain}
\newtheorem{theorem}[subsection]{Theorem}
\newtheorem{proposition}[subsection]{Proposition}
\newtheorem{lemma}[subsection]{Lemma}

\theoremstyle{definition}
\newtheorem{definition}[subsection]{Definition}
\newtheorem{example}[subsection]{Example}
\newtheorem{exercise}[subsection]{Exercise}
\newtheorem{situation}[subsection]{Situation}

\theoremstyle{remark}
\newtheorem{remark}[subsection]{Remark}
\newtheorem{remarks}[subsection]{Remarks}

\numberwithin{equation}{subsection}

% Macros
%
\def\lim{\mathop{\mathrm{lim}}\nolimits}
\def\colim{\mathop{\mathrm{colim}}\nolimits}
\def\Spec{\mathop{\mathrm{Spec}}}
\def\Hom{\mathop{\mathrm{Hom}}\nolimits}
\def\Ext{\mathop{\mathrm{Ext}}\nolimits}
\def\SheafHom{\mathop{\mathcal{H}\!\mathit{om}}\nolimits}
\def\SheafExt{\mathop{\mathcal{E}\!\mathit{xt}}\nolimits}
\def\Sch{\mathit{Sch}}
\def\Mor{\mathop{\mathrm{Mor}}\nolimits}
\def\Ob{\mathop{\mathrm{Ob}}\nolimits}
\def\Sh{\mathop{\mathit{Sh}}\nolimits}
\def\NL{\mathop{N\!L}\nolimits}
\def\CH{\mathop{\mathrm{CH}}\nolimits}
\def\proetale{{pro\text{-}\acute{e}tale}}
\def\etale{{\acute{e}tale}}
\def\QCoh{\mathit{QCoh}}
\def\Ker{\mathop{\mathrm{Ker}}}
\def\Im{\mathop{\mathrm{Im}}}
\def\Coker{\mathop{\mathrm{Coker}}}
\def\Coim{\mathop{\mathrm{Coim}}}

% Boxtimes
%
\DeclareMathSymbol{\boxtimes}{\mathbin}{AMSa}{"02}

%
% Macros for moduli stacks/spaces
%
\def\QCohstack{\mathcal{QC}\!\mathit{oh}}
\def\Cohstack{\mathcal{C}\!\mathit{oh}}
\def\Spacesstack{\mathcal{S}\!\mathit{paces}}
\def\Quotfunctor{\mathrm{Quot}}
\def\Hilbfunctor{\mathrm{Hilb}}
\def\Curvesstack{\mathcal{C}\!\mathit{urves}}
\def\Polarizedstack{\mathcal{P}\!\mathit{olarized}}
\def\Complexesstack{\mathcal{C}\!\mathit{omplexes}}
% \Pic is the operator that assigns to X its picard group, usage \Pic(X)
% \Picardstack_{X/B} denotes the Picard stack of X over B
% \Picardfunctor_{X/B} denotes the Picard functor of X over B
\def\Pic{\mathop{\mathrm{Pic}}\nolimits}
\def\Picardstack{\mathcal{P}\!\mathit{ic}}
\def\Picardfunctor{\mathrm{Pic}}
\def\Deformationcategory{\mathcal{D}\!\mathit{ef}}
\usepackage{cleveref}
\crefname{originallemma}{original lemma}{original lemmas}
\Crefname{originallemma}{Original lemma}{Original lemmas}
\makeatletter
\newlabel{coherent-lemma-cohomology-projective-space-over-ring}{{01XT}{source}{}{sourcepointer-1}{}}
\newlabel{coherent-lemma-cohomology-projective-space-over-ring@cref}{{[originallemma][0][]01XT}{[1][0][]source}{}{}{}}
\makeatother
\crefname{originallemma}{original lemma}{original lemmas}
\Crefname{originallemma}{Original lemma}{Original lemmas}
\makeatletter
\newlabel{coherent-lemma-relative-affine-cohomology}{{089W}{source}{}{sourcepointer-2}{}}
\newlabel{coherent-lemma-relative-affine-cohomology@cref}{{[originallemma][0][]089W}{[1][0][]source}{}{}{}}
\makeatother
\crefname{originallemma}{original lemma}{original lemmas}
\Crefname{originallemma}{Original lemma}{Original lemmas}
\makeatletter
\newlabel{constructions-lemma-proj-sheaves}{{01M7}{source}{}{sourcepointer-3}{}}
\newlabel{constructions-lemma-proj-sheaves@cref}{{[originallemma][0][]01M7}{[1][0][]source}{}{}{}}
\makeatother
\crefname{originalequation}{original equation}{original equations}
\Crefname{originalequation}{Original equation}{Original equations}
\makeatletter
\newlabel{lemma-equation-plane-curve}{{0BYB}{source}{}{sourcepointer-4}{}}
\newlabel{lemma-equation-plane-curve@cref}{{[originalequation][0][]0BYB}{[1][0][]source}{}{}{}}
\makeatother
\title{A plane curve of degree d has trivial global functions and arithmetic genus (d - 1)(d - 2)/2}
\author{The Stacks Project Authors}
\date{stacks-project:curves.tex:local-snapshot-20261012; GNU FDL 1.2 or later}
\begin{document}
\maketitle
\noindent\textit{Editorial note.} All delivered mathematical text is the author's own, copied from the pinned Stacks Project chapter file curves.tex; no statement or proof step was rewritten, reordered or removed.\par\medskip
\noindent\textit{Editorial note.} The counted proof is the author's complete argument: the graded sequence, the sheafified sequence, the cohomology transfer to the plane, and both long-exact identifications with the dimension count are all retained verbatim.\par\medskip
\noindent\textit{Editorial note.} The label lemma-genus-plane-curve names the counted statement, and the proof environment immediately following it in the source is delivered verbatim as its proof.\par\medskip
\noindent\textit{Editorial note.} The counted statement is carried by the author's own lemma environment in the chapter source.\par\medskip
\noindent\textit{Allowed standard prerequisites (editorial).}\begin{itemize}
\item The plane-curve setup. Same chapter, Lemma lemma-equation-plane-curve. Fixes the closed subscheme Z of the projective plane with principal ideal generated by the homogeneous form F of degree d.
\item Graded-to-sheaf exactness and projective-space cohomology. Constructions, Lemma constructions-lemma-proj-sheaves; Cohomology of Schemes, Lemmas coherent-lemma-cohomology-projective-space-over-ring and coherent-lemma-relative-affine-cohomology. Sheafifies 0 -> S(-d) -> S -> S/(F) -> 0 to 0 -> O(-d) -> O -> i\_*O\_Z -> 0, supplies the cohomology of O and O(-d) on the plane, and moves cohomology from Z to the ambient projective space.
\end{itemize}\medskip
\ifdefined\hypertarget\hypertarget{sourcepointer-1}{}\fi
\noindent\textit{Source reference.} Original lemma 01XT: sources/stacks-chow/tags.txt, source lines 2532--2532. Prerequisite pointer invoked at proof line(s) 1929; it names a standard fact of a sibling chapter that the counted argument uses, and every step of the counted proof built on it is retained verbatim in the delivered body.\par\medskip
\ifdefined\hypertarget\hypertarget{sourcepointer-2}{}\fi
\noindent\textit{Source reference.} Original lemma 089W: sources/stacks-chow/tags.txt, source lines 10511--10511. Prerequisite pointer invoked at proof line(s) 1932; it names a standard fact of a sibling chapter that the counted argument uses, and every step of the counted proof built on it is retained verbatim in the delivered body.\par\medskip
\ifdefined\hypertarget\hypertarget{sourcepointer-3}{}\fi
\noindent\textit{Source reference.} Original lemma 01M7: sources/stacks-chow/tags.txt, source lines 2112--2112. Prerequisite pointer invoked at proof line(s) 1919; it names a standard fact of a sibling chapter that the counted argument uses, and every step of the counted proof built on it is retained verbatim in the delivered body.\par\medskip
\ifdefined\hypertarget\hypertarget{sourcepointer-4}{}\fi
\noindent\textit{Source reference.} Original equation 0BYB: sources/stacks-curves/pinned-original-source.tex, source lines 1859--1859. Prerequisite pointer invoked at proof line(s) 1905; it names a standard fact of the same chapter that the counted argument uses, and every step of the counted proof built on it is retained verbatim in the delivered body.\par\medskip
\begin{lemma}
\label{lemma-genus-plane-curve}
Let $Z \subset \mathbf{P}^2_k$ be as in Lemma \ref{lemma-equation-plane-curve}
and let $I(Z) = (F)$ for some $F \in k[T_0, T_1, T_2]$.
Then $H^0(Z, \mathcal{O}_Z) = k$ and the genus of $Z$ is
$(d - 1)(d - 2)/2$ where $d = \deg(F)$.
\end{lemma}

\begin{proof}
Let $S = k[T_0, T_1, T_2]$.
There is an exact sequence of graded modules
$$
0 \to S(-d) \xrightarrow{F} S \to S/(F) \to 0
$$
Denote $i : Z \to \mathbf{P}^2_k$ the given closed immersion.
Applying the exact functor $\widetilde{\ }$
(Constructions, Lemma \ref{constructions-lemma-proj-sheaves})
we obtain
$$
0 \to \mathcal{O}_{\mathbf{P}^2_k}(-d) \to
\mathcal{O}_{\mathbf{P}^2_k} \to i_*\mathcal{O}_Z \to 0
$$
because $F$ generates the ideal of $Z$.
Note that the cohomology groups of $\mathcal{O}_{\mathbf{P}^2_k}(-d)$ and
$\mathcal{O}_{\mathbf{P}^2_k}$ are given in
Cohomology of Schemes, Lemma
\ref{coherent-lemma-cohomology-projective-space-over-ring}.
On the other hand, we have
$H^q(Z, \mathcal{O}_Z) = H^q(\mathbf{P}^2_k, i_*\mathcal{O}_Z)$ by
Cohomology of Schemes, Lemma \ref{coherent-lemma-relative-affine-cohomology}.
Applying the long exact cohomology sequence
we first obtain that
$$
k = H^0(\mathbf{P}^2_k, \mathcal{O}_{\mathbf{P}^2_k}) \longrightarrow
H^0(Z, \mathcal{O}_Z)
$$
is an isomorphism and next that the boundary map
$$
H^1(Z, \mathcal{O}_Z) \longrightarrow
H^2(\mathbf{P}^2_k, \mathcal{O}_{\mathbf{P}^2_k}(-d)) \cong
k[T_0, T_1, T_2]_{d - 3}
$$
is an isomorphism. Since it is easy to see that the dimension of this
is $(d - 1)(d - 2)/2$ the proof is finished.
\end{proof}

\end{document}
